\cleardoublepage
\chapter{ORGANISASI ATAU LINGKUNGAN KERJA PRAKTIK}
\label{chap:organisasi_lingkungan}

\section{Profil Perusahaan}
PT Bank Syariah Indonesia (Persero) Tbk atau BSI resmi beroperasi pada 1 Februari 2021. BSI dibentuk melalui penggabungan PT Bank BRIsyariah Tbk, PT Bank Syariah Mandiri, dan PT Bank BNI Syariah \autocite{bsi2026}. Pada Januari 2026, perseroan memperoleh penyesuaian status administratif sehingga nama perusahaan menjadi PT Bank Syariah Indonesia (Persero) Tbk.

Pembahasan profil perusahaan pada laporan ini dibatasi pada konteks yang berhubungan dengan pekerjaan. Penulis tidak membahas struktur bisnis bank secara keseluruhan karena Kerja Praktik berfokus pada pengolahan data dan analisis untuk campaign.

\section{Posisi Kerja Praktik}
Penulis melaksanakan Kerja Praktik pada fungsi Campaign \& Leads Management di lingkungan Data \& AI Strategy. Fungsi ini bekerja dengan data untuk menyiapkan leads, mendukung pelaksanaan campaign, dan membaca hasil campaign.

Pekerjaan pada fungsi tersebut berhubungan dengan kebutuhan unit bisnis. Kebutuhan dapat berbentuk kriteria calon target, permintaan data, atau kebutuhan analisis. Kriteria tersebut kemudian diterjemahkan menjadi logika data, diperiksa terhadap sumber yang tersedia, dan diproses sebelum digunakan lebih lanjut.

\section{Lingkungan Data dan Tools}
Penulis menggunakan SQL pada lingkungan big data perusahaan melalui antarmuka Hue dengan HiveQL dan Impala. Apache Impala menyediakan antarmuka SQL untuk data pada ekosistem Hadoop dan mendukung operasi seperti \textit{join}, agregasi, \textit{distinct}, \textit{union all}, dan Common Table Expression \autocite{impala2026}. Pada konteks Kerja Praktik, tools tersebut digunakan untuk eksplorasi data, pengembangan query, validasi, serta penyiapan dataset.

Penulis juga menggunakan Talend Data Integration untuk membangun dan menguji alur ETL. Pekerjaan yang telah diuji dapat dijadwalkan melalui Talend Administration Center. Dokumentasi Talend menjelaskan bahwa task dapat dijalankan dengan trigger berbasis waktu maupun event \autocite{talend2026}. Pada pekerjaan penulis, penjadwalan dibahas untuk proses yang berjalan secara berkala.

Spreadsheet digunakan untuk pemeriksaan hasil dan analisis deskriptif. Python digunakan pada pekerjaan spatial clustering dan visualisasi. Pemilihan tools tidak dilakukan berdasarkan satu standar untuk semua pekerjaan. SQL dan Talend lebih sesuai untuk penyiapan data terstruktur dalam alur perusahaan, sedangkan Python lebih sesuai untuk eksperimen analisis dan clustering.

\section{Alur Kerja Data pada Campaign}
Berdasarkan pekerjaan yang diamati penulis, proses dimulai dari kebutuhan unit bisnis. Kebutuhan tersebut diterjemahkan menjadi kriteria data. Penulis kemudian memeriksa ketersediaan atribut, relasi antarsumber, kondisi filter, dan definisi bisnis sebelum membangun query.

Hasil query tidak langsung dianggap benar hanya karena query dapat dijalankan. Penulis memeriksa jumlah baris pada beberapa tahap, duplikasi akibat \textit{join}, nilai kosong, tipe data, serta perubahan jumlah data setelah filter. Pada beberapa kebutuhan, hasil juga melalui proses cleansing dan validasi lain sebelum digunakan sebagai leads.

Alur ini membuat penulis memahami bahwa pekerjaan data di perusahaan mempunyai dua lapisan yang harus berjalan bersama. Lapisan pertama adalah implementasi teknis. Lapisan kedua adalah definisi bisnis. Kesalahan pada salah satu lapisan dapat menghasilkan data yang secara sintaks benar tetapi salah untuk kebutuhan campaign.

\section{Kerahasiaan dan Etika Penggunaan Data}
Data yang digunakan selama Kerja Praktik berada pada lingkungan perusahaan. Sebagian data memuat atribut yang tidak boleh dipindahkan ke media publik. Penulis memisahkan bahan kerja internal dari bahan yang aman untuk laporan akademik.

Pada laporan ini, detail yang dapat mengidentifikasi nasabah atau merekonstruksi struktur data internal dihilangkan. Screenshot internal, query lengkap, nama tabel, identifier, dan nilai bisnis rinci tidak dimuat. Contoh teknis dijelaskan sebagai pola proses. Keputusan ini juga diterapkan pada lampiran.
